% Le forze dei due tendini sulla rotula (problema ud1-vett-rotula), con l'asse x lungo
% la bisettrice: (a) angolo di 120 gradi, R = F; (b) angolo di 90 gradi, R = 1,4 F.
% Scala: 1200 N = 2 cm (R = 2 in (a), R = 2,83 in (b)).
\begin{tikzpicture}[>=Stealth, font=\small]
  \begin{scope}
    \draw[gray, dashed] (-0.3,0) -- (2.6,0) node[right, black] {$x$};
    \coordinate (F1) at (60:2);
    \coordinate (F2) at (-60:2);
    \coordinate (R) at (2,0);
    \draw[dashed, gray] (F1) -- (R) -- (F2);
    \draw[->, very thick, formula] (0,0) -- (F1) node[above] {$\vec F_1$};
    \draw[->, very thick, medica] (0,0) -- (F2) node[below] {$\vec F_2$};
    \draw[->, very thick, esempio] (0,0) -- (R);
    \node[esempio, above] at (1.55,0) {$\vec R$};
    \draw[richiamo, thick] (0.55,0) arc[start angle=0, end angle=60, radius=0.55];
    \draw[richiamo, thick] (0.55,0) arc[start angle=0, end angle=-60, radius=0.55];
    \node[richiamo] at (30:0.9) {\ang{60}};
    \node[richiamo] at (-30:0.9) {\ang{60}};
    \fill (0,0) circle (1.8pt);
    \node[left] at (0,0) {rotula};
    \node at (1,-2.6) {(a) \ang{120}: $R = 2F\cos\ang{60} = F$};
  \end{scope}
  \begin{scope}[xshift=6.2cm]
    \draw[gray, dashed] (-0.3,0) -- (3.3,0) node[right, black] {$x$};
    \coordinate (F1) at (45:2);
    \coordinate (F2) at (-45:2);
    \coordinate (R) at (2.83,0);
    \draw[dashed, gray] (F1) -- (R) -- (F2);
    \draw[->, very thick, formula] (0,0) -- (F1) node[above] {$\vec F_1$};
    \draw[->, very thick, medica] (0,0) -- (F2) node[below] {$\vec F_2$};
    \draw[->, very thick, esempio] (0,0) -- (R);
    \node[esempio, above] at (2.2,0) {$\vec R$};
    \draw[richiamo, thick] (0.55,0) arc[start angle=0, end angle=45, radius=0.55];
    \draw[richiamo, thick] (0.55,0) arc[start angle=0, end angle=-45, radius=0.55];
    \node[richiamo] at (22:0.95) {\ang{45}};
    \node[richiamo] at (-22:0.95) {\ang{45}};
    \fill (0,0) circle (1.8pt);
    \node at (1.4,-2.6) {(b) \ang{90}: $R = 2F\cos\ang{45} \approx 1{,}4\,F$};
  \end{scope}
\end{tikzpicture}
